\documentclass[11pt,a4paper]{article}

\usepackage[utf8]{inputenc}
\usepackage[T1]{fontenc}
\usepackage{amsmath,amssymb}
\usepackage{geometry}
\geometry{margin=2.5cm}

\title{A Minimal Relational Sketch for the Emergence of the Golden Ratio in PDL}
\author{C\'edric Laubscher\\
\small Independent researcher in theoretical and mathematical physics}
\date{February 2026}

\begin{document}
\maketitle

\section{Motivation}

In the PDL framework, composite structures such as the proton are described as hierarchical closures composed of:
\begin{itemize}
  \item an internal quasi-stationary core, characterised by $R_{\text{core}}$ relations (valence content),
  \item a predominantly dynamical background or ``sea'',
  \item an active surface, characterised by $R_{\text{surf}}$ relations, through which coherent coupling to external closures (e.g.\ an electron-type $(4,6)$) is established.
\end{itemize}

The total relational content of the composite is denoted
\[
R_{\text{tot}} = R_{\text{core}} + R_{\text{surf}} + R_{\text{sea}}.
\]

In the proton model used in PDL, the valence content is $r_{\text{valence}} = 930$, the total relational content is $R_{\text{tot}} \approx 11\,017$, and the effective surface is of order $R_{\text{surf}} \approx 5.0 \times 10^2$.

The goal of this note is to sketch how a golden-ratio-like factor can naturally appear when one imposes a minimal optimisation principle relating internal closure and active surface, in a way compatible with the PDL axioms.

\section{Core--Surface Partition and Optimisation}

For present purposes, we focus on the partition between core and active surface and disregard the detailed role of the sea. We write
\[
R_{\text{tot}}' = R_{\text{core}} + R_{\text{surf}},
\]
where $R_{\text{tot}}'$ is the portion of the relational content directly involved in the stationary core and its active surface.

Two antagonistic tendencies are at play:
\begin{itemize}
  \item Internal closure favours a large $R_{\text{core}}/R_{\text{surf}}$ ratio: too small a core weakens the stability of the hierarchical structure.
  \item External coupling favours a large $R_{\text{surf}}/R_{\text{core}}$ ratio: too small a surface restricts the number of coherent interaction channels.
\end{itemize}

A simple way to encode this trade-off is to consider a quality functional of the form
\[
Q(R_{\text{core}},R_{\text{surf}}) = a \,\frac{R_{\text{core}}}{R_{\text{surf}}} + b \,\frac{R_{\text{surf}}}{R_{\text{core}}},
\]
with $a,b>0$. The first term rewards strong internal closure, the second term rewards sufficient external exposure. For fixed $R_{\text{tot}}'$, the ratio
\[
x = \frac{R_{\text{core}}}{R_{\text{surf}}}
\]
is optimised when
\[
\frac{dQ}{dx} = 0 \quad \Rightarrow \quad -a \frac{1}{x^2} + b = 0 \quad \Rightarrow \quad x^2 = \frac{a}{b}.
\]
At this level, the optimum is characterised by $x = \sqrt{a/b}$ and does not yet single out any particular numerical value.

To go beyond this, we note that in PDL the proton is not an isolated object: it belongs to a hierarchy of closures. It is therefore natural to require an additional form of \emph{self-similarity} between the way the total relational budget is partitioned and the way the core and surface themselves relate to each other.

\section{Self-Similar Partition and the Golden Ratio}

We now impose a minimal self-similarity condition on the partition of $R_{\text{tot}}'$:
\begin{equation}
\frac{R_{\text{core}}}{R_{\text{tot}}'} = \frac{R_{\text{surf}}}{R_{\text{core}}}.
\label{eq:self-similar}
\end{equation}
Intuitively, this expresses that:
\begin{itemize}
  \item the role of the core with respect to the whole composite mirrors
  \item the role of the active surface with respect to the core.
\end{itemize}
This is a discrete counterpart of the classical ``part--whole'' characterisations of the golden ratio.

Using $R_{\text{tot}}' = R_{\text{core}} + R_{\text{surf}}$, equation \eqref{eq:self-similar} becomes
\[
\frac{R_{\text{core}}}{R_{\text{core}} + R_{\text{surf}}} = \frac{R_{\text{surf}}}{R_{\text{core}}}.
\]
Let us define the ratio
\[
\lambda = \frac{R_{\text{core}}}{R_{\text{surf}}}.
\]
Then $R_{\text{core}} = \lambda R_{\text{surf}}$ and
\[
\frac{\lambda R_{\text{surf}}}{\lambda R_{\text{surf}} + R_{\text{surf}}} = \frac{R_{\text{surf}}}{\lambda R_{\text{surf}}}
\quad\Rightarrow\quad
\frac{\lambda}{\lambda + 1} = \frac{1}{\lambda}.
\]
This yields the quadratic equation
\[
\lambda^2 = \lambda + 1.
\]
Its positive solution is the golden ratio
\[
\lambda = \varphi = \frac{1 + \sqrt{5}}{2}.
\]
Consequently,
\[
\frac{R_{\text{core}}}{R_{\text{surf}}} = \varphi,\qquad
\frac{R_{\text{core}}}{R_{\text{tot}}'} = \frac{1}{\varphi},\qquad
\frac{R_{\text{surf}}}{R_{\text{tot}}'} = \frac{1}{\varphi^2}.
\]

In this sense, the golden ratio can be seen as the proportion that realises a minimal self-similar partition between core and surface under a simple closure-consistent condition.

\section{Connection with PDL's Proton Architecture}

In the proton model used in PDL, the valence content is $r_{\text{valence}} = 930$, distributed over three quasi-stationary cores (up and down types). A natural way to define an effective core size is
\[
R_{\text{core}} \sim r_{\text{valence}},
\]
while the active surface is taken to be
\[
R_{\text{surf}} \approx \varphi \,\frac{r_{\text{valence}}}{3},
\]
reflecting the fact that the effective surface contribution is distributed over the three valence regimes, modulated by a golden-ratio-like factor.

The above self-similar partition suggests that the factor $\varphi$ should not be regarded as an arbitrary insertion, but as a candidate consequence of:
\begin{itemize}
  \item the requirement of minimal closure for a composite relational structure,
  \item an optimisation mechanism balancing internal coherence and external coupling,
  \item and a self-similarity condition linking the core/whole and surface/core ratios.
\end{itemize}

In the current formulation of PDL, this logic is implemented at the level of $R_{\text{surf}}$ as an informed hypothesis rather than as a fully derived theorem. A more complete treatment would:
\begin{enumerate}
  \item embed the self-similarity condition \eqref{eq:self-similar} into the optimisation functional $\Omega$ used in PDL,
  \item and verify, in a combinatorial setting, that the resulting extremal configurations indeed select golden-ratio-like partitions as the first admissible solutions under the closure constraints.
\end{enumerate}

\section{Outlook}

This note is not intended as a final derivation, but as a minimal sketch showing how the golden ratio can naturally emerge from a discrete optimisation mechanism consistent with the PDL axioms. In this perspective, the appearance of $\varphi$ in the expression
\[
R_{\text{surf}} \approx \varphi \,\frac{r_{\text{valence}}}{3}
\]
should be interpreted as the provisional trace of a deeper closure-oriented optimisation principle, rather than as a numerological insertion.

\end{document}
